\chapter{Diseño e implementación} % Main chapter title

\label{Chapter3} % Change X to a consecutive number; for referencing this chapter elsewhere, use \ref{ChapterX}

Todos los capítulos deben comenzar con un breve párrafo introductorio que indique cuál es el contenido que se encontrará al leerlo.  La redacción sobre el contenido de la memoria debe hacerse en presente y todo lo referido al proyecto en pasado, siempre de modo impersonal.


%----------------------------------------------------------------------------------------
%	SECTION 1
%----------------------------------------------------------------------------------------
\section{Arquitectura de hardware}
Descripción del microcontrolador ESP32 y motivo por el que se seleccionó para este trabajo
\section{Arquitectura de firmware}
Descripción del software desarrollado y módulos principales. Estrategia para garantizar compatibilidad, escalabilidad y mantenibilidad.
\section{Arquitectura de interfaz de usuario}
Descripción de la interfaz de usuario, API implementada, y funcionalidades principales. Estrategia para garantizar compatibilidad, escalabilidad y mantenibilidad
\section{Documentación}
Descripción de la documentación entregada al cliente y obetivo: Documentación del código, guia del desarrollador, guia de usuario y pruebas del sistema



